\chapter{Catálogo de métricas de calidad}
\label{anx:metricas}

Este anexo recoge el catálogo completo de las métricas de calidad de datos
implementadas en \dataqual.  El sistema define \textbf{siete métricas de
calidad}, organizadas en tres categorías según su modo de puntuación, más una
herramienta de diagnóstico no puntuable:

\begin{itemize}
  \item \textbf{Métricas fundamentales} (siempre puntuables): completitud,
        unicidad, exactitud sintáctica y actualidad.
  \item \textbf{Métricas condicionales} (\emph{opt-in}): equilibrio de clases
        y diversidad, que solo contribuyen al \emph{Quality Score} cuando el
        usuario configura explícitamente las columnas a evaluar.
  \item \textbf{Consistencia lógica}: puntuable solo si se definen reglas de
        negocio; devuelve \texttt{score\,=\,None} en caso contrario.
  \item \textbf{Diagnóstico} (no puntuable): detección de valores atípicos,
        disponible exclusivamente desde el módulo de perfilado de datos.
\end{itemize}

Todas las métricas evaluables heredan de la clase abstracta
\texttt{BaseMetric} e implementan el método \texttt{evaluate()}, que recibe
un \texttt{DataFrame} de Pandas y devuelve un objeto \texttt{MetricResult}
con tres campos: \texttt{score} (puntuación $[0,\!1]$ o \texttt{None}),
\texttt{results} (diccionario con datos detallados) e \texttt{issues} (lista
de incidencias detectadas).

\medskip
\noindent\textbf{Cálculo del \emph{Quality Score} global.}\quad
El \emph{Quality Score} \textbf{no} es la media de las puntuaciones
individuales de cada métrica.  En su lugar, parte de $100$\,\% y se reduce
mediante una penalización acumulada basada en el \emph{número y severidad}
de las incidencias detectadas:

\[
  \text{QS} = \max\!\Bigl(0,\;
    1 - \min\!\bigl(0{,}97,\;\tfrac{\text{penalización bruta}}
                                    {\text{factor de escala}}\bigr)\Bigr)
\]

donde la penalización bruta es
$\sum_s n_s \cdot w_s$
con pesos $w_{\text{critical}}=0{,}12$, $w_{\text{high}}=0{,}05$,
$w_{\text{medium}}=0{,}01$ y $w_{\text{low}}=0{,}003$; y el factor de escala
$\sqrt{\max(10,\, C) / 10}$ normaliza por el número de columnas~$C$ del
dataset para que la misma densidad de incidencias produzca la misma
calificación independientemente del ancho del conjunto de datos.


%% =========================================================================
\section{Completitud (\texttt{completeness})}
\label{sec:met-completitud}
%% =========================================================================

\textbf{Objetivo.}\quad
Mide el grado en que los valores esperados están presentes en el conjunto de
datos, alineada con la dimensión de completitud de
\ac{ISO}/\ac{IEC}~25012~\citep{iso25012}.

\textbf{Algoritmo.}\quad
Para cada columna se calcula la ratio de valores no nulos:
$\text{completitud}_c = 1 - \overline{\text{nulos}}_c$.
La puntuación global es la media aritmética de las ratios de las columnas
seleccionadas (o de todas, si no se especifican).

Antes del cálculo, se aplican los \textbf{patrones de nulidad}
configurados (\texttt{null\_patterns}): cadenas vacías, literales como
\texttt{"null"}, \texttt{"N/A"}, \texttt{"nan"}, \texttt{"none"}, guiones
y espacios en blanco se convierten a \texttt{NaN} para que no inflen
artificialmente la completitud.  El usuario puede definir patrones
\emph{preset} (7 predefinidos) o expresiones regulares personalizadas.

\textbf{Generación de incidencias.}\quad
Se emiten dos tipos de \emph{issues}:
\begin{enumerate}
  \item \textbf{Global}: si la completitud media del dataset cae por debajo
        del umbral configurado.
  \item \textbf{Por columna}: cada columna individual cuya completitud
        esté por debajo del umbral genera su propia incidencia.
\end{enumerate}

\textbf{Parámetros de configuración:}
\begin{table}[htbp]
\centering\small
\begin{tabularx}{\linewidth}{llcX}
\toprule
\tabhead{Parámetro} & \tabhead{Tipo} & \tabhead{Defecto} & \tabhead{Descripción} \\
\midrule
\texttt{columns}        & lista   & todas & Columnas a evaluar. \\
\texttt{threshold}      & float   & 0,95  & Umbral mínimo de completitud. \\
\texttt{null\_patterns}  & dict    & ---   & \texttt{\{presets:~[...], custom:~[...]\}}: patrones de nulidad adicionales. \\
\bottomrule
\end{tabularx}
\end{table}


%% =========================================================================
\section{Unicidad (\texttt{uniqueness})}
\label{sec:met-unicidad}
%% =========================================================================

\textbf{Objetivo.}\quad
Evalúa la ausencia de registros duplicados a nivel de fila y la variabilidad
por columna individual.

\textbf{Algoritmo.}\quad
La evaluación opera en dos niveles simultáneos:

\begin{enumerate}
  \item \textbf{Unicidad de filas}: $U_{\text{filas}} =
        \text{filas\_únicas} / \text{total\_filas}$.
  \item \textbf{Variabilidad por columna}: para cada columna se calcula
        $V_c = n_{\text{únicos}}^{(c)} / n_{\text{no nulos}}^{(c)}$
        y se compara con un \textbf{umbral adaptativo} que depende del tipo
        inferido de la columna:
        \begin{itemize}
          \item Columnas de identificador (\texttt{*\_id}, \texttt{*\_uuid},
                \texttt{*\_key}): umbral = 0,95.
          \item Columnas categóricas ($\leq 20$ valores únicos, no numéricas):
                umbral = 0,05.
          \item Resto de columnas: umbral = 0,30.
        \end{itemize}
\end{enumerate}

La \textbf{puntuación combinada} pondera ambos niveles cuando existen
columnas de identificador:
$S = 0{,}7 \cdot U_{\text{filas}} + 0{,}3 \cdot \overline{V_{\text{ID}}}$;
si no las hay, $S = U_{\text{filas}}$.

\textbf{Generación de incidencias.}\quad
Se emiten tres tipos de \emph{issues}:
\begin{enumerate}
  \item \textbf{Baja variabilidad}: columnas cuya variabilidad está por debajo
        de su umbral adaptativo.  La severidad para columnas no-ID se calcula
        con distancia relativa al umbral (no absoluta) para evitar falsos
        críticos en columnas categóricas legítimas.
  \item \textbf{Filas duplicadas}: si $U_{\text{filas}} < \text{threshold}$, se
        adjuntan hasta 5 filas de ejemplo (con enmascaramiento de columnas
        sensibles).
  \item \textbf{Identificador no único}: columnas que el usuario marca como
        únicas pero contienen valores repetidos.
\end{enumerate}

\textbf{Parámetros de configuración:}
\begin{table}[htbp]
\centering\small
\begin{tabularx}{\linewidth}{llcX}
\toprule
\tabhead{Parámetro} & \tabhead{Tipo} & \tabhead{Defecto} & \tabhead{Descripción} \\
\midrule
\texttt{threshold} & float & 1,0  & Umbral mínimo de unicidad de filas. \\
\texttt{columns}   & lista & ---  & Columnas que deben ser estrictamente únicas. \\
\bottomrule
\end{tabularx}
\end{table}


%% =========================================================================
\section{Exactitud sintáctica (\texttt{syntactic\_accuracy})}
\label{sec:met-exactitud}
%% =========================================================================

\textbf{Objetivo.}\quad
Verifica que los valores de cada columna se ajustan a patrones de formato
predefinidos (correo electrónico, teléfono, \ac{DNI}, etc.) o identificadores
únicos globales (\ac{UUID}), alineada con la dimensión de exactitud de
\ac{ISO}/\ac{IEC}~25012~\citep{iso25012}.

\textbf{Catálogo de patrones integrados.}\quad
El sistema incluye 13~expresiones regulares predefinidas:

\begin{table}[htbp]
\centering\small
\begin{tabularx}{\linewidth}{lX}
\toprule
\tabhead{Tipo} & \tabhead{Ejemplo / Descripción} \\
\midrule
\texttt{email}           & \texttt{user@example.com} \\
\texttt{url}             & \texttt{https://...} \\
\texttt{phone\_es}       & Teléfono español: \texttt{+34 6XXXXXXXX} \\
\texttt{phone\_intl}     & Teléfono internacional genérico \\
\texttt{dni\_es}          & \ac{DNI} español: 8 dígitos + letra \\
\texttt{date\_iso}        & Fecha ISO: \texttt{AAAA-MM-DD} \\
\texttt{date\_eu}         & Fecha europea: \texttt{DD/MM/AAAA} \\
\texttt{integer}         & Número entero (con signo) \\
\texttt{decimal}         & Número decimal (coma o punto) \\
\texttt{uuid}            & Identificador \ac{UUID} (cualquier versión) \\
\texttt{ip\_v4}           & Dirección IPv4 \\
\texttt{postal\_code\_es} & Código postal español (5 dígitos) \\
\texttt{credit\_card}    & Tarjeta de crédito (4 grupos de 4 dígitos) \\
\bottomrule
\end{tabularx}
\end{table}

Además, los usuarios pueden definir \textbf{patrones personalizados}
almacenados en la base de datos (\texttt{ValidationPattern}), que se cargan
en tiempo de ejecución y sobrescriben los predefinidos cuando comparten la
misma clave.

\textbf{Auto-detección de tipos.}\quad
Para columnas de tipo \texttt{object} sin configuración explícita, se
muestrea hasta 100~valores y se prueban todos los patrones del catálogo.
Se asigna el tipo con mayor tasa de coincidencia \textbf{si y solo si}
supera el 85\,\%.  Si dos o más tipos superan el umbral simultáneamente,
la columna se marca como \textbf{formato mixto} (\texttt{mixed\_format})
y se emite un \emph{issue} de severidad media en lugar de asignar
arbitrariamente un tipo.

\textbf{Parámetros de configuración:}
\begin{table}[htbp]
\centering\small
\begin{tabularx}{\linewidth}{llcX}
\toprule
\tabhead{Parámetro} & \tabhead{Tipo} & \tabhead{Defecto} & \tabhead{Descripción} \\
\midrule
\texttt{columns}           & lista & --- & Lista de \texttt{\{column, expected\_type, pattern\}}. \\
\texttt{auto\_detect\_types} & bool  & true & Activar auto-detección para columnas sin tipo. \\
\texttt{custom\_patterns}   & dict  & \{\} & Mapa \texttt{\{columna: regex\}} adicional. \\
\texttt{threshold}         & float & 0,95 & Umbral mínimo de conformidad. \\
\bottomrule
\end{tabularx}
\end{table}


%% =========================================================================
\section{Consistencia lógica (\texttt{logical\_consistency})}
\label{sec:met-consistencia}
%% =========================================================================

\textbf{Objetivo.}\quad
Valida reglas de negocio entre columnas del dataset, expresadas como
expresiones de consulta o reglas condicionales \texttt{IF--THEN}.

\textbf{Tipos de reglas soportados:}
\begin{enumerate}
  \item \textbf{Violación directa} (\texttt{type:~"violation"}):
        la expresión selecciona las filas que violan la regla.  Tasa de
        cumplimiento:
        $1 - \text{violaciones} / \text{total\_filas}$.
  \item \textbf{Condicional} (\texttt{type:~"if\_then"}):
        se evalúa \texttt{condition} para seleccionar el ámbito y después
        se verifica \texttt{assertion} sobre ese subconjunto.  La tasa de
        cumplimiento se normaliza sobre las filas que cumplen la condición,
        no sobre el total del dataset, para que una regla selectiva con
        100\,\% de violaciones en su ámbito se refleje como un problema
        grave.
\end{enumerate}

Las reglas también se pueden expresar en lenguaje natural con la sintaxis
\texttt{IF <condición> THEN <aserción>} (o sus equivalentes en español
\texttt{Si... Entonces}), que el motor parsea automáticamente mediante
expresión regular.

\textbf{Seguridad.}\quad
Antes de ejecutar cualquier regla, el sistema verifica que no contenga
\emph{tokens} prohibidos (\texttt{import}, \texttt{exec}, \texttt{eval},
\texttt{os.}, \texttt{subprocess}, \texttt{lambda}, etc.) para prevenir
inyección de código.  Las reglas bloqueadas generan un \emph{issue} de
severidad alta.

Si no se configuran reglas, la métrica devuelve \texttt{score\,=\,None}
y no contribuye al \emph{Quality Score}.

\textbf{Parámetros de configuración:}
\begin{table}[htbp]
\centering\small
\begin{tabularx}{\linewidth}{llcX}
\toprule
\tabhead{Parámetro} & \tabhead{Tipo} & \tabhead{Defecto} & \tabhead{Descripción} \\
\midrule
\texttt{rules}  & lista & --- & Lista de reglas: \texttt{\{name, type, expression, condition, assertion\}}. \\
\bottomrule
\end{tabularx}
\end{table}


%% =========================================================================
\section{Equilibrio de clases (\texttt{class\_balance})}
\label{sec:met-balance}
%% =========================================================================

\textbf{Objetivo.}\quad
Cuantifica el equilibrio de la distribución de variables categóricas,
especialmente relevante para la variable objetivo en tareas de \ac{ML}
supervisado.

\textbf{Algoritmo.}\quad
Se utiliza la entropía de Shannon~(\citeyear{shannon1948}):
\[
  H = -\sum_{i} p_i \cdot \log_2(p_i), \qquad
  H_{\max} = \log_2(n_{\text{clases}}), \qquad
  \text{Índice de equilibrio} = \frac{H}{H_{\max}} \times 100
\]

La métrica distingue dos modos de operación:
\begin{itemize}
  \item \textbf{Modo explícito} (columnas configuradas por el usuario):
        la puntuación contribuye al \emph{Quality Score}.
  \item \textbf{Modo auto-detectado}: se analizan automáticamente las
        columnas categóricas con hasta 50~valores únicos (y ratio de
        unicidad $\leq 25\,$\% para excluir identificadores), pero sus
        resultados son solo informativos y \textbf{no} penalizan la
        calificación global, ya que el desequilibrio natural (p.~ej.,
        97/3 en detección de fraude) no es un defecto de calidad per~se
        según \ac{ISO}/\ac{IEC}~5259~\citep{iso5259}.
\end{itemize}

\textbf{Conformidad con distribución esperada.}\quad
Opcionalmente, el usuario puede definir rangos porcentuales esperados para
cada clase (p.~ej., \texttt{junior:~[70,~85]}, \texttt{senior:~[10,~20]}).
Se calcula una puntuación de conformidad basada en cuántas clases caen dentro
de su rango, y la puntuación final combina entropía y conformidad al 50\,\%.

\textbf{Parámetros de configuración:}
\begin{table}[htbp]
\centering\small
\begin{tabularx}{\linewidth}{llcX}
\toprule
\tabhead{Parámetro} & \tabhead{Tipo} & \tabhead{Defecto} & \tabhead{Descripción} \\
\midrule
\texttt{columns}             & lista & ---  & Columnas objetivo (modo explícito). \\
\texttt{auto\_detect}         & bool  & true & Activar auto-detección. \\
\texttt{max\_cardinality}     & int   & 50   & Máx.\ valores únicos para auto-detección. \\
\texttt{imbalance\_threshold\_high} & float & 0,90 & Umbral de clase dominante. \\
\texttt{imbalance\_threshold\_low}  & float & 0,05 & Umbral de clase minoritaria. \\
\texttt{expected\_distribution} & dict & \{\} & Rangos esperados por clase: \texttt{\{clase: [min\%, max\%]\}}. \\
\bottomrule
\end{tabularx}
\end{table}


%% =========================================================================
\section{Actualidad (\texttt{currentness})}
\label{sec:met-actualidad}
%% =========================================================================

\textbf{Objetivo.}\quad
Mide la frescura de los datos calculando la antigüedad del registro de fecha
más reciente de cada columna temporal.  Implementa la dimensión
\emph{Currentness} ($\Delta T_2$) de \ac{ISO}/\ac{IEC}~5259-2:2024 §6.2.5, medida
Cur-ML-1 (\emph{Feature currentness}).  No debe confundirse con
\emph{Timeliness} ($\Delta T_1$), que mide la latencia de ingesta y requiere
\emph{timestamps} externos.

\textbf{Algoritmo de frescura.}\quad
La puntuación de frescura por columna sigue una degradación lineal con
factor de escala configurable (\texttt{DECAY\_SCALE\,=\,3}):

\[
  \text{frescura} =
  \begin{cases}
    1{,}0 & \text{si } d \leq U \\[4pt]
    \max\!\bigl(0,\; 1 - \frac{d - U}{U \cdot k}\bigr)
      & \text{si } d > U
  \end{cases}
\]

donde $d$ son los días de antigüedad, $U$ el umbral de obsolescencia
y $k = 3$ el factor de escala.  Con el umbral por defecto de 30~días,
la puntuación alcanza 0 a los 120~días.

\textbf{Auto-detección.}\quad
Columnas con \texttt{dtype=datetime64} se incluyen directamente.
Para columnas de texto, se parsea una muestra de 50~valores y se incluyen
si al menos el 50\,\% se reconoce como fecha.  Si la tasa de parseo de una
columna incluida es inferior al 80\,\%, se genera un \emph{issue}
informativo adicional.

\textbf{Severidad de obsolescencia.}\quad
La severidad se determina por la ratio $d / U$:
\emph{medium} ($\geq 1\times$), \emph{high} ($\geq 3\times$) y
\emph{critical} ($\geq 10\times$).

\textbf{Parámetros de configuración:}
\begin{table}[htbp]
\centering\small
\begin{tabularx}{\linewidth}{llcX}
\toprule
\tabhead{Parámetro} & \tabhead{Tipo} & \tabhead{Defecto} & \tabhead{Descripción} \\
\midrule
\texttt{columns}                  & lista & ---  & Columnas temporales a evaluar. \\
\texttt{auto\_detect}              & bool  & true & Activar auto-detección. \\
\texttt{staleness\_threshold\_days} & int   & 30   & Umbral de obsolescencia en días. \\
\bottomrule
\end{tabularx}
\end{table}


%% =========================================================================
\section{Diversidad (\texttt{diversity})}
\label{sec:met-diversidad}
%% =========================================================================

\textbf{Objetivo.}\quad
Mide si el dataset cubre adecuadamente los valores o rangos que el usuario
espera encontrar en cada columna, alineada con la dimensión de diversidad
de \ac{ISO}/\ac{IEC}~5259-2~\citep{iso5259}.

Esta métrica es \textbf{opt-in por columna}: solo se evalúan las columnas
para las que el usuario define expectativas.  Si no hay expectativas
definidas, devuelve \texttt{score\,=\,None} y queda excluida del
\emph{Quality Score}.

\textbf{Evaluación categórica.}\quad
Para columnas de tipo \texttt{categorical}, se calcula la ratio de
valores esperados presentes en el dataset.  Cada valor presente con
representación adecuada ($\geq$ \texttt{min\_representation}) recibe
crédito completo; los presentes pero subrepresentados reciben medio
crédito; los ausentes, cero.

\textbf{Evaluación numérica.}\quad
Para columnas de tipo \texttt{numeric}, se divide el rango esperado
$[\min, \max]$ en $n$~bins uniformes y se calcula la cobertura como
la ratio de bins con al menos un dato:
$S = \text{bins\_con\_datos} / \text{total\_bins}$.

\textbf{Parámetros de configuración:}
\begin{table}[htbp]
\centering\small
\begin{tabularx}{\linewidth}{llcX}
\toprule
\tabhead{Parámetro} & \tabhead{Tipo} & \tabhead{Defecto} & \tabhead{Descripción} \\
\midrule
\texttt{columns}            & dict  & \{\} & Mapa de columnas con sus expectativas. \\
\texttt{threshold}          & float & 0,60 & Umbral mínimo de diversidad. \\
\bottomrule
\end{tabularx}
\end{table}

Para cada columna el usuario especifica el tipo (\texttt{categorical}
o \texttt{numeric}) junto con \texttt{expected\_values} (lista de valores
esperados) o \texttt{expected\_range} ($[\min, \max]$) y opcionalmente
\texttt{min\_representation} (defecto:~1\,\%) o \texttt{num\_bins}
(defecto:~10).


%% =========================================================================
\section{Valores atípicos (\texttt{outliers})}
\label{sec:met-outliers}
%% =========================================================================

\textbf{Nota.}\quad
Esta clase \textbf{no es una métrica evaluable}: no produce puntuación
ni contribuye al \emph{Quality Score}.  No está registrada en el
\texttt{METRIC\_REGISTRY} y se invoca exclusivamente desde el módulo
de perfilado de datos (\emph{Data Profiling}), en línea con la
recomendación de \ac{ISO}/\ac{IEC}~5259~\citep{iso5259}.

\textbf{Métodos de detección.}\quad
El sistema soporta dos métodos:

\begin{enumerate}
  \item \textbf{IQR} (por defecto): un valor~$x$ es atípico si
    $x < Q_1 - f \cdot \mathrm{IQR}$ o
    $x > Q_3 + f \cdot \mathrm{IQR}$, donde $f$ es el factor
    configurable.
  \item \textbf{Z-score}: un valor es atípico si
    $|z| > f$, donde $z = (x - \mu) / \sigma$.
\end{enumerate}

Cada incidencia reporta la proporción de valores atípicos, hasta
5~valores de ejemplo (enmascarados si la columna es sensible) y las
estadísticas de la distribución ($Q_1$, $Q_3$, \ac{IQR} o $\mu$, $\sigma$).

\textbf{Parámetros de configuración:}
\begin{table}[htbp]
\centering\small
\begin{tabularx}{\linewidth}{llcX}
\toprule
\tabhead{Parámetro} & \tabhead{Tipo} & \tabhead{Defecto} & \tabhead{Descripción} \\
\midrule
\texttt{columns} & lista  & todas las numéricas & Columnas a evaluar. \\
\texttt{method}  & string & \texttt{iqr}        & Método: \texttt{iqr} o \texttt{zscore}. \\
\texttt{factor}  & float  & 1,5                 & Factor multiplicador ($f$). \\
\bottomrule
\end{tabularx}
\end{table}


%% =========================================================================
\section{Modelo de severidad dinámica}
\label{sec:met-severidad}
%% =========================================================================

Las métricas asignan a cada incidencia una \textbf{severidad dinámica} en
función de la distancia entre el valor real y el umbral configurado.  Los
criterios se definen en \texttt{BaseMetric.calculate\_dynamic\_severity()} y
dependen del tipo de incidencia, no de la métrica concreta: las cuatro
dimensiones en las que «mayor es mejor» (completitud, unicidad, exactitud
sintáctica y consistencia lógica) comparten un mismo criterio, mientras que
equilibrio de clases, actualidad y diversidad emplean criterios propios.  El
Cuadro~\ref{tab:severidad} los resume.

\begin{table}[htbp]
\centering\small
\caption{Criterios de severidad dinámica por tipo de incidencia.}
\label{tab:severidad}
\begin{tabularx}{\linewidth}{@{}>{\raggedright\arraybackslash}p{3.7cm}XXXX@{}}
\toprule
\tabhead{Tipo de incidencia} & \tabhead{Critical} & \tabhead{High} & \tabhead{Medium} & \tabhead{Low} \\
\midrule
Completitud, unicidad, exactitud sintáctica, consistencia lógica & $< 50\,$\% & $< 70\,$\% o dist.\,$> 15$\,pp & dist.\,$> 5$\,pp & dist.\,$\leq 5$\,pp \\
Equilibrio: clase dominante & $\geq 99\,$\% & $\geq 95\,$\% & $\geq 90\,$\% & $< 90\,$\% \\
Equilibrio: desviación de distribución esperada & $\geq 20$\,pp & $\geq 10$\,pp & $\geq 5$\,pp & $< 5$\,pp \\
Actualidad (ratio $d/U$) & $\geq 10\times$ & $\geq 3\times$ & $\geq 1\times$ & $< 1\times$ \\
Diversidad & $\leq 20\,$\% & $\leq 40\,$\% & $<$ umbral & $\geq$ umbral \\
Outliers (proporción)\,$^{*}$ & $\geq 20\,$\% & $\geq 10\,$\% & $\geq 5\,$\% & $< 5\,$\% \\
\bottomrule
\end{tabularx}
\par\smallskip
{\footnotesize $^{*}$\,Outliers es una herramienta de \emph{profiling}: no
contribuye al \emph{Quality Score} ni figura en el registro de métricas, pero
sus incidencias reciben severidad con el mismo modelo.}
\end{table}


% Local Variables:
%  coding: utf-8
%  mode: latex
%  mode: flyspell
%  ispell-local-dictionary: "castellano8"
% End:
